% 译者说明（Lead）
\chapter*{译者说明}
\addcontentsline{toc}{chapter}{译者说明}
\markboth{译者说明}{译者说明}

本译本依据河野俊丈《場の理論とトポロジー》（岩波书店，1998 年）译出。原书是作者在岩波讲座基础上写成的单行本，正文三章，另有卷首的「理论概要と展望」一篇导论。

\medskip
\noindent\textbf{体例说明}
\begin{itemize}[leftmargin=2em,itemsep=2pt]
\item 原书页码以 \verb|% [原文 p.N]| 的形式标在每页起始处；PDF 页号 = 原书页码 + 20。卷首导论用罗马数字页码。
\item 公式编号、定义/定理/命题等陈述编号一律照录原书，不重排、不改号。
\item 原书正文的笔误照原样排，只在边注中（\verb|% [Tutor Note: ...]|）指出并给出正确形式；批注统一收在书末的《译注汇总》里。
\item 术语以书末《术语对照表》为准；人名、模型名、缩写保留英文原形。
\item 图均从原书扫描页裁出，图号与图注照原书。
\item 每节前有「本节重点」，每章末有「本章回顾与评注」，均为译者所加，用以提示结构与易错点，可在导言区用 \verb|\showtakeawayfalse|、\verb|\showchapterreviewfalse| 关闭。
\item 参考文献照录原书（不翻译）；原书索引未收，见 README 的「已知限制」。
\end{itemize}

\medskip
\noindent 译文力求准确，但错误在所难免，恳请读者指正。
